\documentclass{article}
\usepackage{graphicx} % Required for inserting images

\title{Шпора МАТСТАТ}
\author{Аким Тарасов}


\usepackage[T2A]{fontenc}
\usepackage[utf8]{inputenc}
\usepackage[russian]{babel}
\usepackage{amsmath}
\usepackage{amssymb}
\usepackage{amsthm}
\usepackage{mathtools}

\date{May 2026}

\begin{document}

\section*{1) ЗАДАЧА НА ОЦЕНКИ (Т3-Т6 из задания)}

Пусть дана случайная величина и выборка
\[
\xi \sim \rho(x,\theta), \vec{x}_n, 
\]
где \(\theta \in \Theta\) --- неизвестный параметр. Требуется построить оценку параметра, исследовать её свойства, проверить эффективность и построить доверительные интервалы.

\subsection*{1. Построение оценки}

\subsubsection*{1.1. Метод моментов (ОММ)}

\begin{enumerate}
    \item Выписать теоретические моменты:
    \[
    \alpha_k(\theta)=M x^k.
    \]

    \item Выписать выборочные моменты:
    \[
    \widetilde{\alpha}_k=\frac1n\sum_{i=1}^n x_i^k.
    \]

    \item Приравнять:
    \[
    \alpha_k(\theta)=\widetilde{\alpha}_k.
    \]

    \item Решить полученное уравнение относительно \(\theta\):
    \[
    \widetilde{\theta}_{OMM}=g(\widetilde{\alpha}_1,\dots,\widetilde{\alpha}_m).
    \]
\end{enumerate}

\subsubsection*{1.2. Метод максимального правдоподобия (ОМП)}

\begin{enumerate}
    \item Выписать функцию правдоподобия:
    \[
    L(\theta)=\prod_{i=1}^n \rho(x_i, \theta).
    \]

    \item Перейти к логарифмической функции правдоподобия:
    \[
    \ell(\theta)=\ln L(\theta,\vec{x}_n).
    \]

    \item Найти стационарные точки:
    \[
    \frac{\partial \ell(\theta)}{\partial \theta}=0.
    \]

    \item Учесть ограничения на параметр:
    \[
    \theta \in \Theta.
    \]

    Особенно важно проверить:
    \begin{itemize}
        \item границы области допустимых значений;
        \item случаи, когда носитель распределения зависит от параметра.
    \end{itemize}

    \item Найти
    \[
    \widetilde{\theta}_{OMP}
    =
    \arg\max_{\theta \in \Theta} L(\theta,X).
    \]
\end{enumerate}

\subsection*{2. Исследование оценки на несмещённость}

Оценка называется несмещённой, если
\[
M\widetilde{\theta}=\theta.
\]

\begin{enumerate}
    \item Вычислить:
    \[
    M \widetilde{\theta}.
    \]

    \item Если
    \[
    M \widetilde{\theta}\neq \theta,
    \]
    то оценка смещённая.

    \item При необходимости построить исправленную оценку:
    \[
    \widetilde{\theta}^{\,*}=a\widetilde{\theta}+b,
    \]
    где коэффициенты \(a,b\) выбираются из условия
    \[
    M\widetilde{\theta}^{\,*}=\theta.
    \]
\end{enumerate}

\subsection*{3. Исследование оценки на состоятельность}

Оценка называется состоятельной, если
\[
\widetilde{\theta}\xrightarrow[]{P}\theta,
\qquad n\to\infty.
\]

\subsubsection*{3.1. Достаточное условие состоятельности}

Если оценка несмещена, или
\[
M\widetilde{\theta}=\theta
\]
и
\[
D\widetilde{\theta}\to0,
\qquad n\to\infty,
\]
то по неравенству Чебышёва
\[
\widetilde{\theta}\xrightarrow[]{P}\theta.
\]

\subsubsection*{3.2. Проверка по определению}

Если достаточное условие не выполняется, исследовать:
\[
P_\theta\left(
|\widetilde{\theta}-\theta|>\varepsilon
\right),
\]
и показать, что
\[
P\left(
|\widetilde{\theta}-\theta|>\varepsilon
\right)\to0
\qquad \forall \varepsilon>0.
\]

\subsection*{4. Исследование эффективности оценки}

\subsubsection*{4.1. Проверка условий регулярности модели}

Для применения неравенства Крамера--Рао требуется регулярность вероятностной модели:

\begin{enumerate}
    \item Существует производная
    \[
    \frac{\partial}{\partial\theta}
    \ln \rho(x, \theta).
    \]

    \item Допустима перестановка производной по параметру и интеграла Стилтьеса:
    \[
    \frac{\partial}{\partial\theta}
    \int
    \,dF(x,\theta)
    =
    \int
    \,d\!\left(
    \frac{\partial F(x,\theta)}{\partial\theta}
    \right).
    \]

    На практике данное условие обычно проверяется следующим образом:
    \[
    \int
    \frac{\partial}{\partial\theta}
    p(x, \theta)\,dx
    +
    \sum_k
    \frac{\partial}{\partial\theta}
    p_k(\theta)
    =
    0,
    \]
    где:
    \begin{itemize}
        \item \(p(x, \theta)\) --- непрерывная часть плотности;
        \item \(p_k(\theta)\) --- вероятности атомов дискретной части распределения.
    \end{itemize}

    \item Информация Фишера
    \[
    I(\theta)
    =
    M[
    \left(
    \frac{\partial}{\partial\theta}
    \ln \rho(x, \theta)
    \right)^2]
    \]
    существует, непрерывна и положительна:
    \[
    I(\theta)>0, \forall \theta \in \Theta
    \]
\end{enumerate}

\subsubsection*{4.2. Проверка регулярности оценки}

Требуется:
\begin{itemize}
    \item несмещённость оценки;
    \item регулярность модели;
    \item ограниченность дисперсии на любом компакте из \[\Theta\].
\end{itemize}

\subsubsection*{4.3. Неравенство Крамера--Рао}

Для любой регулярной несмещённой оценки:
\[
D\widetilde{\theta}
\ge
\frac1{nI(\theta)}.
\]

\begin{enumerate}
    \item Вычислить:
    \[
    D \widetilde{\theta}.
    \]

    \item Вычислить информацию Фишера:
    \[
    I(\theta).
    \]

    \item Сравнить:
    \[
    D \widetilde{\theta}
    \quad \text{и} \quad
    \frac1{nI(\theta)}.
    \]
\end{enumerate}

\subsubsection*{4.4. Вывод}

\begin{itemize}
    \item Если
    \[
    D \widetilde{\theta}
    =
    \frac1{nI(\theta)},
    \]
    то оценка эффективна.

    \item Если
    \[
    D\widetilde{\theta}
    >
    \frac1{nI(\theta)},
    \]
    то не можем ничего сказать, вспоминаем про Теорему о единственности эффектичной оценки.
\end{itemize}

\subsection*{5. Построение доверительных интервалов}

\subsubsection*{5.1. Точный доверительный интервал}

Если удаётся построить статистику
\[
f(\vec{x}_n, h).
\]
распределение которой не зависит от параметра \(\theta\), то используется точный метод.

\begin{enumerate}
    \item Найти распределение статистики:
    \[
    f(\vec{x}_n, h) \sim q(t)
    \]

    \item Построить вероятностное утверждение:
    \[
    P\left(
    q_{(1-\beta)/2}<f(\vec{x}_n, h)<q_{(1+\beta)/2}
    \right)=\beta.
    \]

    \item Разрешить неравенство относительно \(\theta\).
\end{enumerate}

\subsubsection*{5.2. Асимптотический доверительный интервал для ОММ}

Пусть
\[
\widetilde{\theta}_{OMM}
=
g(\widetilde{\alpha}_1,\dots,\widetilde{\alpha}_m),
\]
где \(g\) дифференцируема и
\[
\nabla g(\alpha)\neq0.
\]

Тогда:
\[
\sqrt n
\left(
\widetilde{\theta}_{OMM}-\theta
\right)
\Rightarrow
N(0,\sigma^2(\theta)),
\]
где
\[
\sigma^2(\theta)
=
\nabla^T g(\alpha)\,
K\,
\nabla g(\alpha),
\]
а \(K\) --- ковариационная матрица моментов.

По лемме Слуцкого:
\[
\sigma^2(\theta)
\sim
\sigma^2(\widetilde{\theta}).
\]

Следовательно:
\[
\frac{
\sqrt n
(\widetilde{\theta}_{OMM}-\theta)
}{
\sigma(\widetilde{\theta}_{OMM})
}
\Rightarrow
N(0,1).
\]

Отсюда асимптотический доверительный интервал:
\[
\widetilde{\theta}_{OMM}
\pm
\frac{
z_{1-\alpha/2}\,
\sigma(\widetilde{\theta}_{OMM})
}{
\sqrt n
}.
\]

\subsubsection*{5.3. Асимптотический доверительный интервал для ОМП}

Если модель сильно регулярна (регулярная, перестановочна вторая производная и F 2-жды непрерывно дифференцируема) и
\[
\Theta
\]
--- открытое множество, то:
\[
\sqrt n
(\widetilde{\theta}_{OMP}-\theta)
\Rightarrow
N\left(
0,
I^{-1}(\theta)
\right).
\]

Если рассматривается функция параметра:
\[
g(\theta),
\]
\[
\sqrt n
\left(
g(\widetilde{\theta}_{OMP})-g(\theta)
\right)
\Rightarrow
N\left(
0,
\nabla^T g(\theta)\,
I^{-1}(\theta)\,
\nabla g(\theta)
\right).
\]

После применения леммы Слуцкого:
\[
\frac{
\sqrt n
(g(\widetilde{\theta}_{OMP})-g(\theta))
}{
\sigma(\widetilde{\theta}_{OMP})
}
\Rightarrow
N(0,1).
\]

Итоговый асимптотический доверительный интервал:
\[
g(\widetilde{\theta}_{OMP})
\pm
\frac{
z_{1-\alpha/2}\,
\sigma(\widetilde{\theta}_{OMP})
}{
\sqrt n
}.
\]





\section*{2) ЗАДАЧА НА НАИБОЛЕЕ МОЩНЫЙ КРИТЕРИЙ (ТЕОРЕМА НЕЙМАНА--ПИРСОНА) (Т11 из задания)}

Пусть дана выборка
\[
\vec x_n,
\]
и требуется проверить простые гипотезы:
\[
H_0 : p_0(x),
\qquad
H_1 : p_1(x).
\]

Требуется построить наиболее мощный критерий уровня значимости \(\alpha\), найти критическую область, ошибку первого рода, ошибку второго рода и мощность критерия.

\subsection*{1. Построение отношения правдоподобия}

Выписать функции правдоподобия:
\[
L_0(\vec x_n),
\qquad
L_1(\vec x_n).
\]

Построить отношение правдоподобия:
\[
l(\vec x_n)
=
\frac{L_1(\vec x_n)}
{L_0(\vec x_n)}.
\]

Упростить выражение и исследовать его монотонность.

Обычно удаётся свести критерий к некоторой статистике:
\[
\sum X_i,
\qquad
\overline X,
\qquad
X_{(1)},
\qquad
X_{(n)}
\]
и т.д.

\subsection*{2. Построение критической области}

По теореме Неймана--Пирсона наиболее мощный критерий имеет критическую область:
\[
G
=
\{
\vec x_n :
l(\vec x_n)>C
\}.
\]

После преобразования обычно получается область вида:
\[
G
=
\{
T(\vec X_n)>c
\}
\]
или
\[
G
=
\{
T(\vec X_n)<c
\}.
\]

\subsection*{3. Определение константы из уровня значимости}

Константа \(c\) определяется из условия:
\[
P(G|H_0)=\alpha.
\]

Для вычисления вероятности используются:
\begin{itemize}
    \item известные распределения статистик;
    \item теорема Фишера;
    \item интегрирование плотности;
    \item функции распределения.
\end{itemize}

После нахождения \(c\) критическая область считается полностью построенной.

\subsection*{4. Ошибка первого рода}

Ошибка первого рода:
\[
\alpha_1
=
P(G|H_0).
\]

Это вероятность отвергнуть \(H_0\), когда гипотеза \(H_0\) верна.

\subsection*{5. Мощность критерия}

Мощность критерия определяется как:
\[
W
=
P(G|H_1).
\]

То есть вероятность попадания в критическую область при справедливости альтернативной гипотезы.

Часто мощность выражается через интеграл:
\[
W
=
\int_G p_1(x)\,dx.
\]

\subsection*{6. Ошибка второго рода}

Ошибка второго рода:
\[
\alpha_2
=
1-W.
\]

Это вероятность принять \(H_0\), когда верна гипотеза \(H_1\).

\subsection*{7. Состоятельность критерия}

Критерий называется состоятельным, если:
\[
\forall H_1 \qquad => \qquad
W_n\to1,
\qquad
n\to\infty.
\]

Обычно исследуется предел мощности:
\[
W_n
=
\int_{a_n}^{+\infty} f(x)\,dx
\]
или
\[
W_n
=
\int_{-\infty}^{b_n} f(x)\,dx.
\]

Если при \(n\to\infty\):
\[
a_n\to-\infty
\]
или
\[
b_n\to+\infty,
\]
то интеграл берётся по всей области распределения, следовательно:
\[
W_n\to1.
\]

Тогда:
\[
\alpha_{2}^n\to0,
\]
и критерий является состоятельным.

\end{document}
